\documentclass[10pt,a4paper]{article}
\usepackage[T1]{fontenc}
\usepackage{lmodern,amsmath,amssymb,microtype,xcolor,geometry,hyperref}
\geometry{margin=19mm}
\definecolor{softgray}{gray}{0.94}
\hypersetup{hidelinks,hypertexnames=false,pdftitle={SGA 2 Expose VIII Proposition 3.2 proof opening},pdfauthor={Interlanguage project}}
\newcommand{\cF}{\mathcal F}
\newcommand{\cG}{\mathcal G}
\newcommand{\cO}{\mathcal O}
\newcommand{\prof}{\operatorname{prof}}
\newcommand{\codimn}{\operatorname{codim}}
\begin{document}
\begin{center}
  {\Large\bfseries SGA 2: Local Cohomology and Lefschetz Theorems}\par\smallskip
  {\large Expos\'e VIII: The Finiteness Theorem}\par
  \small Bounded source-aligned working unit --- 19 July 2026
\end{center}
\noindent\colorbox{softgray}{\parbox{0.95\linewidth}{\footnotesize\textbf{Authority.}
Corrected French lines 2901--2907; original printed p.~97; physical
source-PDF pp.~84--85; recomposed running pp.~76--77. The unit crosses a
recomposed page boundary, but the next \texttt{pageoriginale} marker is only
at French line 2915, so its original printed-page coordinate remains p.~97.
Blank line 2908 is excluded. French line 2909, the implication
\textup{(a)}$\Rightarrow$\textup{(b)}, is the exact continuation cursor. The
compiled French reader is generated from the same corrected edition and is
page/layout evidence, not independent corroboration. The symbol $Y$ is
inherited from Corollary~2.3 and the Section~2 convention $U=X\setminus Y$;
this contextual clarification is not a French-source emendation. Under closed
Option~A, target calligraphic $\cF,\cG$ are established English
normalizations from source upright $F,G$, not literal glyph preservation.}}

\medskip
\noindent\textit{Proof, opening.}---Let $\cG$ be a coherent $\cO_X$-module
such that $\cG|_U\simeq\cF$ (EGA~I, 9.4.3). Applying Corollary~2.3 to
$\cG$, we find that condition \textup{(a)} is equivalent to \textup{(c)}:
\begin{enumerate}
\item[(c)] for every $x\in S$, one has
\[
  \prof\cF_x>n-c(x),
  \qquad
  c(x)=\codimn(\overline{\{x\}}\cap Y,\overline{\{x\}}).
\]
\end{enumerate}
\end{document}
